% !PRINTED 240
\weberpage{240}
%<<CONT
arises through composition of $D_2$ with
\[
j_1 = \begin{pmatrix} 0, & -\sqrt{i} \\ i\sqrt{i}, & 0 \end{pmatrix}
\]
and thus yields
\[
D_2' = D_2 + D_2 j_1\text{.}
\]
Now in general:
\[
j_1^{-1} \begin{pmatrix} \alpha, & \beta \\ \gamma, & \delta \end{pmatrix} j_1 = \begin{pmatrix} \delta, & \gamma i \\ -\beta i, & \alpha \end{pmatrix},
\]
and from this one concludes by \S.\,56, (4), that $j_1^{-1} T j_1 = T$, and hence that from the tetrahedral group two extended groups $T' = T+Tj$ and $T'' = T+Tj_1$ can be derived; but also only two, because by \S.\,56 the tetrahedral group is completely determined by a four-group contained in it. For if, instead of the extension $D_2'$, we choose the extension of the four-group to one of the groups (11), we obtain an extended tetrahedral group $T'''$, which passes into $T''$ by transformation with $\varphi_1$ or $\varphi_2$.

Since the substitution
\[
\begin{pmatrix} 0, & e^{\frac{\pi i}{8}} \\ e^{-\frac{\pi i}{8}}, & 0 \end{pmatrix}
\]
does not commute with the octahedral group \S.\,57, (8), no further group of the second kind can be derived from the octahedral group except by inversion, and the same holds of the icosahedral group. To apply this mode of reasoning, one must transform the icosahedral group in such a way that one of the four-groups contained in it takes the simplest form, which is somewhat lengthy, but by no means difficult, and shall not be carried out further here.

We thus obtain the following further groups of the second kind:
\[
\begin{aligned}
&\text{IV.}\quad 1)\ T' = T+Tj, \qquad 2)\ T'' = T+Tj_1.\\
&\text{V.}\qquad\qquad O' = O+Oj.\\
&\text{IV.}\qquad\qquad J' = J+Jj.
\end{aligned}
\]
% [sic] the third line is printed "IV." again (as the first); expected "VI." given the sequence — reproduced as printed.
